\documentclass[a4wide]{article}

\title{Rebuttal}

\input{config}
\author{{}}




\begin{document}

\maketitle

%====================================
\section*{General response}

We thank all reviewers for their careful reading and constructive feedback. We are glad that they find the main theoretical results correct, the exposition clear, the polyhedral results and their connections with the analysis of FW methods novel and useful.

In particular, regarding significance and novelty, our main goal is to both shed light on the conditioning of product polytopes, and to understand how their geometry affects the convergence rates of FW-type algorithms for convex functions satisfying the PL inequality.
In this context, we provide the first explicit composition formulas for two key constants describing that geometry, that are essential in establishing linear convergence of FW methods on polytopes. In doing so, we close a meaningful gap in the FW literature that it is technically nontrivial, and had not been addressed before. 
Several large-scale applications routinely require optimization over product-structured polyhedral constraints, and practitioners already use FW variants in this setting as LMOs decompose naturally across blocks. Having explicit, interpretable condition numbers for these domains tells us when fast FW variants, such as the away-step one, can be expected to behave well. In this direction, we quantify how adding blocks (i.e., additional polyhedral constraints) to the product degrades performance, namely how the rates deteriorate as a function of the conditioning of the individual blocks.
Furthermore, we view our formulas for the two important polytope conditioning numbers, and associated rates, as a first step toward such a systematic theory and as a tool for analyzing existing FW variants and developing new ones in block-separable polyhedral settings.

Regarding our computational experiments, the concern raised by one of the reviewers is that our current log-log plots of gap versus wall-clock time do not make the linear regime visually obvious in some instances. We acknowledge this presentation issue and are rerunning the experiments with a larger iteration budget. 
In the final version, if the paper is accepted, we will replace the current plots with these runs and redact Section 6 to enhance its clarity and improve its structure.

We also welcome the constructive feedback on other parts of the manuscript, and will make sure to integrate the reviewers' suggestions whenever possible in the final version.







%====================================
\section*{Reviewer wei5}

\subsection*{Review}
Summary And Contributions:
This paper addresses the convergence behavior of Frank–Wolfe algorithms over structured constraint sets. The authors analyze the Away Frank–Wolfe (AFW) method on product polytopes, establishing linear convergence rates for smooth convex functions that satisfy the Polyak–Łojasiewicz (PL) condition. The work introduces new geometric quantities—the pyramidal width and the vertex–facet distance—specifically characterized for product polytopes, and derives explicit composition rules showing how these condition numbers behave under Cartesian products. Leveraging these results, the paper demonstrates that AFW achieves linear convergence over product polytopes and applies the theory to an approximate feasibility problem for intersections of multiple polytopes. Empirical experiments further support the theoretical findings, confirming the predicted convergence behavior in high-dimensional settings.

Paper Keywords: Frank-wolfe algorithm, product polytopes, linear convergence
Expertise Keywords: ML optimization theory, learning theory, mathematical optimization

Soundness: Correct / minor errors (e.g., typo or errors that do not affect the main results)
Soundness Justification:
The theoretical analysis appears rigorous and internally consistent. The main theorems on the composition of condition numbers for product polytopes (Theorem 4.5 and Proposition 4.7) are clearly derived and supported by complete proofs. The convergence results for the AFW algorithm under the PL condition follow standard arguments in the Frank–Wolfe literature and are correctly specialized to the product setting. The experimental section is consistent with the theoretical predictions, showing empirical linear convergence. I did not find any major logical or mathematical flaws beyond minor presentation issues.

Significance: Somewhat significant (e.g., theoretical contribution of limited novelty or empirical gains missing baselines/statistical rigor).
Significance Justification:
\textcolor{red}{The theoretical results are well-developed, but their overall significance is moderate. The extension of Frank–Wolfe convergence analysis to product polytopes is technically solid yet conceptually close to existing results for single polytopes. While the derived composition rules for pyramidal width and vertex–facet distance are interesting from a geometric standpoint, their broader impact appears to be limited, as product polytopes represent a relatively straightforward, structured domain. The application to the polytope intersection feasibility problem is illustrative but not clearly more effective than classical projection-based approaches.}

Novelty: New results
Novelty Justification:
The paper presents new theoretical results by deriving explicit composition rules for the pyramidal width and vertex–facet distance over product polytopes. These findings appear original and are not immediately from prior work on single polytopes. \textcolor{red}{Nonetheless, the core idea is a natural extension of established Frank–Wolfe analyses, and thus the novelty lies mainly in the technical refinement rather than a conceptual breakthrough.}

Non-conventional Contributions:
The paper contains some non-conventional theoretical elements. For example, the introduction of the affine pyramidal width as an equivalent but more tractable variant of the standard pyramidal width provides a novel analytical tool that simplifies the derivation of geometric condition numbers for product polytopes. This reformulation is conceptually interesting and may prove useful beyond the specific setting of this paper. While the overall framework follows established Frank–Wolfe theory, this specific construction represents an original and creative methodological contribution.

Clarity:
Except for the points I raised above, this paper is well-written and easy to follow.

Relation To Prior Work: All related works are clearly discussed
Additional Comments:
N/A

Reproducibility: Sufficient amount of details available for reproducing the main results

Rating: 5: Borderline Accept (Technically solid paper where reasons to accept outweigh reasons to reject, e.g., limited evaluation.)
Confidence: 3: Fairly confident (It is possible that you did not understand some parts of the submission or that you are unfamiliar with some pieces of related work. Math/other details were not carefully checked.)

\subsection*{Response}

We thank the reviewer for carefully reading the paper and for the positive assessment on the correctness and clarity of our results.

Regarding the significance of the product setting and the degree of novelty, we would like to clarify our perspective. Many practical applications across several fields involve Cartesian products of polyhedral constraints, where in particular the FW LMO decomposes naturally and can be called independently on each block. 
While this may appear to be a ``natural'' extension of the single-polytope setting, existing FW analyses do not provide explicit formulas for the pyramidal width or vertex–facet distance of product domains. As a consequence, prior work does not quantify how these quantities scale as the number of polytopes in the product grows, nor how one should expect FW convergence rates to deteriorate as a function of this scaling.
Notably, our affine pyramidal width is designed precisely to make this structure more easily analyzable, as it is equivalent to the classical pyramidal width on general polytopes but it admits exact algebraic composition rules on products.
We view this as an important step toward a systematic understanding of FW-type methods on block-structured domains with multiple polytope constraints. We already take a first step in this direction, by quantifying how the AFW rate depends on the number and conditioning of the individual blocks, and by studying an application to approximate feasibility over intersections of several polytopes.

Regarding this application, we chose it precisely because it is a classical and actively studied problem in convex optimization, with a substantial literature on alternating projection methods (in particular for the two-set case). In this context, our aim is not to argue that AFW universally outperforms projection-based schemes, but rather to show that AFW offers a practically relevant alternative whose linear convergence on polytopes can be quantified via our product-polytope condition numbers. We actually did implement an alternating projections baseline in our computational experiments. However, on our instances the per-iteration cost of computing projections onto each polytope (even approximately) grew so quickly with the number of sets and the size of the polytopes that it became prohibitive. As a result, running alternating projections for a comparable number of iterations was too expensive, and a fair empirical comparison was far from feasible within our experimental budget.










%====================================
\section*{Reviewer 3Yb1}

\subsection*{Review}
Summary And Contributions:
This paper considers the question of "how linear convergence results for Frank-Wolfe-type algorithms extend from single polytopes to product polytopes". The paper studies two geometric condition numbers, the pyramidal width and the vertex-facet distance, showing how these quantities behave under Cartesian products. Given the numbers for individual polytopes, the authors provide lower bounds on the numbers for the product polytope. Based on these results, the paper analyzes the Away-step Frank-Wolfe (AFW) algorithm and shows that it guarantees linear convergence under convexity, smoothness, and the PL condition.

Paper Keywords: Away-step Frank-Wolfe, linear convergence, product polytope
Expertise Keywords: Frank-Wolfe, linear convergence, polytope

Soundness: Correct / minor errors (e.g., typo or errors that do not affect the main results)
Soundness Justification: N/A

Significance: Somewhat significant (e.g., theoretical contribution of limited novelty or empirical gains missing baselines/statistical rigor).
Significance Justification:
This paper establishes rigorous theoretical contributions toward understanding Frank-Wolfe-type algorithms. It provides lower bounds for the pyramidal width and vertex-facet distance of the Cartesian product of multiple polytopes based on those of the individual polytopes. Using this characterization of the condition numbers, the paper provides linear convergence results for AFW under convexity, smoothness, and the PL condition.

Novelty: Incremental compared to existing results
Novelty Justification:
The main contributions of this paper are on characterizing the two condition number measures for product polytopes. Proofs are based on convex and polyhedral analysis techniques. Upon establishing the lower bounds on the condition numbers, linear convergence is expected.

Non-conventional Contributions:
No non-conventional contributions

Clarity:
The paper is very well-written.

Relation To Prior Work: All related works are clearly discussed

Reproducibility: Sufficient amount of details available for reproducing the main results

Rating: 5: Borderline Accept (Technically solid paper where reasons to accept outweigh reasons to reject, e.g., limited evaluation.)
Confidence: 3: Fairly confident (It is possible that you did not understand some parts of the submission or that you are unfamiliar with some pieces of related work. Math/other details were not carefully checked.)

\subsection*{Response}

We thank the reviewer for the clear summary and the positive evaluation of the technical correctness and writing.

We would like to clarify where we see the main novelty. Our aim is to make the conditioning of product polytopes tractable and to link this with the beahavior of FW-type methods. On the geometric side, we (i) show that the pyramidal width of the Cartesian product of two polytopes admits an \emph{exact} closed-form expression in terms of the widths of the factors (Theorem 4.5), and (ii) for a general product of $k>2$ polytopes we derive an exact formula for its vertex-facet distance and a lower bound on its pyramidal width in terms of the individual widths that is tight (Corollary 4.4, Proposition 4.7, Corollary A.3).
Previous FW analyses only treat single polytopes and do not provide such composition rules for a generic number $k>1$ of factors. Building on these geometric results, we can quantify linear convergence rates for AFW on polytope product domains (under convexity, smoothness and a PL condition), which depend explicitly on the resulting condition numbers.
Further, beyond their geometric interest, we believe that these formulas sharpen the understanding of the pyramidal width and vertex–facet distance as key geometric constants underpinning existing linear convergence analyses of FW variants, and may serve as a starting point for further work on FW methods in structured domains.










%====================================
\section*{Reviewer FpBL}

\subsection*{Review}
Summary And Contributions:
This paper investigates the linear convergence properties of Away-step Frank-Wolfe (AFW) algorithms over product domains of polytopes, a setting that naturally arises in multi-block convex optimization problems. The authors introduce the notion of affine pyramidal width, that is a geometric measure capturing the conditioning of product polytopes. They also derive exact composition formulas relating the width of a product polytope to that of its component polytopes. These geometric results are used to establish a linear convergence rate for AFW under a PL condition on the objective. The theoretical results are supported by synthetic and feasibility-based experiments.

Paper Keywords: Conditional Gradient Methods, Convergence Analysis, PL-condition
Expertise Keywords: Frank-Wolfe style Algorithm, Convergence Analysis, Convex Optimization

Soundness: Correct / minor errors (e.g., typo or errors that do not affect the main results)
Soundness Justification:
The main geometric result (Theorem 4.5) correctly establishes that the affine pyramidal width of a product polytope satisfies $\delta_{\mathcal{P} \times \mathcal{Q}} = \dfrac{\sqrt{\delta_{\mathcal{P}}^2 \delta_{\mathcal{Q}}^2 }}{\sqrt{\delta_{\mathcal{P}^2} + \delta_{\mathcal{Q}^2}}}$.
The proof is logically coherent and uses standard convex geometry (support functions, face decompositions, and projection orthogonality) without hidden assumptions. The dependence on the product's ambient space dimension is properly normalized, preventing scaling artifacts. The step transitioning from Lemma 4.2 (face decomposition lemma) to Lemma 4.3 (affine hull equivalence) is valid and non-trivial-the orthogonal projection argument avoids an earlier mistake found in Pena \& Rodriguez (2018) where the affine hull was implicitly assumed full-dimensional. This is an excellent correction and is technically sound.
PL-based linear convergence proof: The convergence analysis in Section 4.4 builds on the standard FW template: upper bounding the dual gap via the pyramidal width and linking it to the decrease in objective under a $\mu$-PL condition. Proposition 4.8 correctly applies the generalized Descent Lemma under L-smoothness and uses the standard quadratic decrease argument. The proof structure is tight; each inequality corresponds to a defined geometric constant.

Significance: Somewhat significant (e.g., theoretical contribution of limited novelty or empirical gains missing baselines/statistical rigor).
Significance Justification:
The work closes a significant gap in the literature: while Pena \& Rodríguez (2018) and Garber \& Hazan (2016) characterized pyramidal widths and linear rates for single polytopes, no prior work derived explicit composition formulas for product domains. Product structures like multi-block optimization, separable constraints, and distributed convex feasibility are pervasive that fit this template. This extension is therefore meaningful and timely.
The results justify why AFW retains linear convergence on high-dimensional, block-separable domains. This fact previously observed empirically in exsiting literature, but not explained theoretically. This has immediate implications for practical problems such as matrix factorization and block-sparse regularization, where LMOs operate independently on each block.
The coupling between the PL condition and pyramidal width in this paper offers a geometric analogue of the condition number used in gradient descent analysis. This strengthens the conceptual connection between conditional gradient methods and first-order smooth optimization.

Novelty: New results
Novelty Justification:
The paper introduces the affine pyramidal width, rigorously proving its equivalence to the classical pyramidal width for single polytopes but demonstrating its superior analytical behavior under product composition. This is genuinely novel.
By working under the PL condition rather than full strong convexity, the paper extends the known applicability of linear convergence proofs. This theoretical relaxation is important because it covers many structured convex losses used in ML, e.g., logistic regression and squared hinge losses.
Unlike earlier works that derived only asymptotic bounds ($\delta_{\mathcal{P} \times \mathcal{Q}} \ge  \min\{\delta_{\mathcal{P}}, \delta_{\mathcal{P}}\} / \sqrt{2}$), this paper provides an exact algebraic relationship between component widths. This analytical precision is rare and represents a methodological advance.

Non-conventional Contributions:
The geometric insights in this paper on affine pyramidal width are nontrivial and adds a new lens to conditioning analysis.

Clarity:
The paper is well-structured: Section 2 lays out background, Section 3 introduces definitions, and Section 4 delivers the main geometric and convergence results. The logic progresses naturally.

Relation To Prior Work: All related works are clearly discussed
Additional Comments: \textcolor{red}{The constants in proposition 4.9's convergence rate (the 32 factor and the denominator in  $\rho = \mu \delta^2 / 32 L D^2$) appear without derivation. While this doesn't undermine the asymptotic claim, the absence of clear constant tracking weakens reproducibility. I could not confirm whether 32 arises from a single use of Cauchy-Schwarz or a cumulative bound across blocks.
Several transitions, for e.g., between Lemma 4.3 and Theorem 4.5 is abrupt, with insufficient narrative linking the geometric argument to its analytical role in convergence proofs. A short “concept bridge” paragraph summarizing the intuition behind the affine hull projections would improve readability.
The influence of the number of blocks (k) on the convergence rate is only discussed qualitatively (“minor polynomial factor away”). A clear derivation of how $\delta_{\mathcal{P} \times \mathcal{Q}}$  scales with k would improve the completeness of the result.}

Reproducibility: Sufficient amount of details available for reproducing the main results

Rating: 6: Accept (Technically solid paper, with high impact on at least one sub-area. The contribution is convincing and the evaluation is adequate, though there may be some limitations or open questions.)
Confidence: 2: Somewhat confident (You are willing to defend your assessment, but it is quite likely that you did not understand central parts of the submission or that you are unfamiliar with some pieces of related work. Math/other details were not carefully checked.)


\subsection*{Response}

We thank the reviewer for the careful reading and the detailed feedback.

The rates in our Proposition~4.9 are obtained by instantiating Theorem~3.1 from ``Linearly Convergent Away-Step Conditional Gradient for Non-strongly Convex Functions'' (by Beck and Shtern, 2015) with the parameters of our setting, and then translating their geometric decay into an explicit bound on iterations $t$.
In particular, this last step is carried out exactly as in the proof of our Proposition 4.8. Concretely, the geometric decay of Theorem~3.1 is
\[
f(x_k) - f^\ast \le C\bigl(1 - \alpha^\dagger\bigr)^{(k-1)/2}\,,
\]
where
\[
\alpha^\dagger
=
\min\left\{
  \frac{\Omega_X^2}{8 \rho \kappa D^2 N^2},
  \frac{1}{2}
\right\}\,.
\]
In our notation, $\Omega_X$ is the vertex–facet distance $\alpha_{\mathsf{vf}}$ of the domain polytope $\mathcal{P}$; $\rho$ is the smoothness constant that we call $L$; $\kappa$ in our context is defined as $\kappa = 2/\mu$, where $\mu$ is the quadratic-growth/PL parameter (see Definition~3.2); $D$ is the set diameter, that we denote by $D_{\mathcal{P}}$, $N$ is the upper bound on the active-set size discussed in Appendix~B (notably, either $|V_P|$ or $nk+1$ for the product polytope). Finally, we set $C = f(x_t) - f(x^\ast)$. Plugging these into $\alpha^\dagger$ yields
\[
\alpha^\dagger
=
\min\left\{
\frac{\mu\alpha_{\mathsf{vf}}^2}{16LD_{\mathcal{P}}^2N^2},
\frac{1}{2}
\right\}.
\]
As in Proposition~4.8, we then use $(1-x)^r \le \exp(-rx)$ for $r>0$, $x\in(0,1)$ to obtain
\[
f(x_t) - f(x^\ast) \le
\Bigl(f(x_t) - f(x^\ast)\Bigr) \exp\Bigl(- \frac{\mu\mathrm{vf}_P^2}{32 L D_{\mathcal{P}}^2 N^2} (t-1)\Bigr)\,.
\]
Solving $f(x_t) - f(x^\ast) \le \varepsilon$ for $t$, with $\epsilon$ a given optimality tolerance, gives precisely the iteration complexity stated in Proposition 4.9. In the camera-ready version, we will add a short paragraph in the main text explicitly referring to Theorem~3.1 of Beck and Shtern and summarizing this ``constant tracking'', so that the origin of the factor $32$ and the other terms is transparent.

We agree that the transition from Lemma~4.3 to Theorem~4.5 is conceptually important and deserves a clearer bridge in the main text. Lemma~4.3 shows that the point realizing the affine pyramidal width lies in the relative interior of the corresponding face, which immediately implies Corollary~4.4. This equivalence is what allows us to work with the affine pyramidal width in the proof of Theorem~4.5 while still obtaining the classical pyramidal width for the product. We will add a brief sentence before Theorem~4.5 emphasizing that Lemma~4.3 states a key structural property making the affine variant both equivalent to the classical pyramidal width, and better suited for our compositional arguments.

Lastly, the reviewer asks for a more quantitative explanation of our statement that the lower bound on the product pyramidal width is only a ``minor polynomial factor away'' from the trivial upper bound.
This is exactly what Theorem 4.5 and Corollary A.3 examine, as already stated right after Theorem~4.5 in the main text. 
Together with the upper bound $\delta_{\mathcal{P}} \le \min_i \delta_{\mathcal{P}_i}$, they show that our bounds are tight up to a constant $\sqrt{k}$: in other words, the pyramidal width of the product can deteriorate by at most a factor $\sqrt{k}$.
In the revision, we will state this bound directly and make the dependence on $k$ fully explicit from the outset.











%====================================
\section*{Reviewer 8ayF}

\subsection*{Review}
Summary And Contributions:
This paper studies the rate of linear convergence with Frank-Wolfe algorithms on a convex set comprised of the product of polytopes. On well-behaved functions such as strong convex and (slightly generalized) Polyak-Lojasiewicz functions, we can enjoy linear convergence with Away Frank-Wolfe (AFW), but its convergence rate is affected by the nature of the domain. Specifically, when the domain is a polytope, the convergence rate is affected by its pyramidal width and vertex-facet distance. This paper determines these condition numbers of the product of polytopes when every polytope's condition number is known, leading to derive the convergence rate of the Frank-Wolfe on the product of polytopes. An application of this finding to the problem of approximately finding a point in the intersection of multiple polytopes is empirically exhibited.

Paper Keywords: Frank-Wolfe, pyramidal width, vertex-facet distance
Expertise Keywords: Frank-Wolfe, convex optimization, potential game

Soundness: \textcolor{red}{Major errors} (e.g., an incorrect theorem or derivation)
Soundness Justification:
The main technical contribution of this paper is to derive the pyramidal width and the vertex-facet distance of the product of polytopes (Theorem 4.5 and Proposition 4.7). As far as I read, their proofs, described in Appendices A.1 and A.2, are correct. I didn't check the proofs of other claims (like the convergence rate of AFW) in Appendix, but I think that these results are correct because they are largely a simple combination of the main results (Theorem 4.5 and Proposition 4.7) and the existing results.
\textcolor{red}{However, the linear convergence is not well justified by the experiments. In the figures, time vs the primal/FW gap was plotted with log-log plot. If we enjoy linear convergence, i.e. the gap decays exponentially w.r.t. iterations, the plot will be concave down. while we observe such a trend in Figure 1, in Figures 2 and 3 we cannot see such a trend. On log-log plot, linear trend means that the gap decays polynomially w.r.t. iterations. Probably I suspect that the number of iterations are insufficient to observe linear convergence for these experiments; the primal/FW gaps stayed fairly large (~1e3,1e4) for these experiments.}

Significance: Significant contributions (e.g., strong theoretical insights or well-supported empirical improvements with appropriate baselines/statistics).
Significance Justification:
I think the technical contribution to derive the condition numbers of the product of polytope is significant because the convex optimization on the product of polytope frequently appears in applications, like exhibited in the paper. I also suggest that such a problem also arises in the computation of equilibrium in a (continuous) potential games, like selfish routing, with multiple commodities. Deriving a condition number of such a domain is important to assess the convergence rate of the Frank-Wolfe variant algorithms.
Although the linear convergence is not well confirmed by the experiments (see the Soundness section), I believe that these technical contributions are still significant.

Novelty: New results
Novelty Justification:
The main technical contribution is novel; as written in Introduction, while a convex optimization on the product of polytopes frequently appears in a number of applications, the convergence of Frank-Wolfe variants on such a problem has not been known.

Non-conventional Contributions:
N/A

Clarity:
The technical parts, Sections 3 to 5, are well structured and organized; it is easy to follow the technical details of claims. Compared to this, \textcolor{red}{the description for experiments (Section 6) are highly non-structured. The prerequisites for the experiments are described in the latter part of Section 6, and the descriptions for experimental results are scattered over the text. Authors should restructure this part to separate the prerequisites, the results, and the discussions/observations clearly. Moreover, as described in Soundness section, I think that more iteration will be needed to observe the linear convergence for the experiments.}

Relation To Prior Work: All related works are clearly discussed
Additional Comments:
\textcolor{red}{I believe that Figures 1--3 are log-log plots because the horizontal tics are written as $10^x$, but it is unclear. Please specify somewhere that these are log-log plots.}

Reproducibility: Sufficient amount of details available for reproducing the main results

Rating: 4: Weak Accept / Weak Reject (Paper has some merits but also notable weaknesses. Reasons to accept and reject are roughly balanced, and the paper may require further evaluation or discussion. Please use sparingly.)
Confidence: 3: Fairly confident (It is possible that you did not understand some parts of the submission or that you are unfamiliar with some pieces of related work. Math/other details were not carefully checked.)


\subsection*{Response}

We thank the reviewer for the detailed and insightful comments, and for highlighting both the polyhedral contributions and their relevance for first-order convex optimization methods.

We would first like to clarify that the plots in Section 6 and Appendix C display the primal and FW gaps versus \emph{elapsed wall-clock time} (in seconds), not versus iteration count, and confirm that both axes use a log–log scale (in base $\log_{10}$).
We acknowledge that, in this format, the linear regime for AFW is indeed a bit harder to see, especially when the runs are stopped while the curves are still in a relatively high-gap regime, as is the case for the most part in Figures 2 and 3. We are therefore currently rerunning the experiments with a substantially larger iteration budget to drive the gaps much lower and to better expose the geometric decrease in gaps. In the camera-ready version, if the paper is accepted, we will replace the current plots with new ones after the currently running experiments are finished. We expect that these longer runs will also address the reviewer's concern about the relatively large gaps observed at the end of some of the current runs.

Regarding the computational test we carried out, we would like to clarify that our main theoretical results (Theorem 4.5, Proposition 4.7, Proposition 4.9, and their corollaries) are correct (as the reviewer remarks) and, in that sense, independent from the experiments in Section 6. The concern raised is that the current plots do not make the linear regime visually obvious in some instances. While we agree that this is a limitation of how the experiments are currently presented and we are working to fix it, we view it as an issue of empirical illustration and exposition rather than a soundness error in the theoretical results themselves.

Regarding the numerical experiments, we would like to emphasize that our main theoretical results (Theorem 4.5, Proposition 4.7, Proposition 4.9 and their corollaries) are proved independently of Section 6 and, as the reviewer notes, their proofs appear correct. The concern raised is that the current plots do not make the linear regime visually obvious in some instances. We agree that
this is a limitation of how the experiments are currently presented and are working to improve it. However, we view this as an issue of empirical illustration and exposition, rather than a problem with the soundness of the theoretical results themselves. We hope that this clarifies why we do not believe the paper should be regarded as having ``major errors''.

We appreciate the reviewer's suggestion regarding the structure of the experimental section. We will more carefully reorganize it to have a cleaner narrative flow and clearly separate: (i) a setup subsection, describing in one place algorithmic variants, hyperparameters and stopping criteria, the instance generation protocols (for intersecting and disjoint polytopes), and software/hardware environment; (ii) a results subsection, devoted exclusively to the qualitative and quantitative assessment of the observed performances. We will also make explicit at the beginning of Section 6 that all figures plot $\log_{10}$-gaps versus wall-clock time.

We are grateful for the reviewer's suggestion that similar product-polytope structures arise in the computation of equilibria in continuous potential games. We agree that this is an interesting direction and that our composition formulas for pyramidal width and vertex–facet distance could help quantify convergence rates of FW-type dynamics in such settings. We are not aware of a reference that explicitly formulates these equilibrium problems in the exact product polytope framework we study. We would be thankful if the reviewer could point us to a specific source, that we would be happy to cite in the final version.




%====================================
% \bibliographystyle{plainnat}
% \bibliography{bib.bib}

\clearpage
\appendix
\onecolumn
\end{document}
